% =============================================================
%  MAT231-F27 Calculus II — Lecture 09
%  Unit 9: R^n: norms, open and closed sets, limits and continuity
%
%  Build from the repository root:   npm run materials -- lecture-09
%  This writes lecture-09.pdf (slides, with overlays) and
%  lecture-09-handout.pdf (print version) into this folder.
% =============================================================
\documentclass{wosm-slides}

\usepackage{pgfplots}
\pgfplotsset{compat=1.18}

\newcommand{\R}{\mathbb{R}}
\newcommand{\norm}[1]{\lVert #1 \rVert}
\newcommand{\vx}{\mathbf{x}}
\newcommand{\va}{\mathbf{a}}
\newcommand{\vy}{\mathbf{y}}

\course{MAT231-F27}{Calculus II}
\title[L09 \textperiodcentered\ Functions of several variables]{Functions of several variables}
\subtitle{$\mathbb{R}^n$: norms, open and closed sets, limits and continuity}
\author{Dr.\ Ali Khudiyev}
\institute{WoSM Academy}
\date{Lecture 09\enspace\textperiodcentered\enspace Unit 9\enspace\textperiodcentered\enspace Fall 2027}

\begin{document}

\maketitle

% -------------------------------------------------------------
\begin{frame}{Where we are}
  \begin{itemize}
    \item \textbf{Units 1--8}: series, power series, Taylor series, generating functions,
      Stirling's formula --- all for functions of \emph{one} variable.
    \item \textbf{Units 9--17}: differential calculus in $\R^n$. Partial derivatives, the gradient,
      the chain rule and the Jacobian, the Hessian, optimisation, gradient descent, Lagrange
      multipliers.
    \item Then multiple integrals (Units 18--22) and differential equations (Units 23--25).
  \end{itemize}
  \vskip0.6em
  Today we set the stage: distance in $\R^n$, which sets are open or closed, and what a limit
  means when you can approach a point from infinitely many directions.
  Vectors and inner products come from Linear Algebra (MAT102-S27).
\end{frame}

\begin{frame}{This lecture}
  \tableofcontents
  \vskip1em
  {\small\color{wosm-muted}
  Reading: G.~Strang, E.~Herman, \emph{Calculus, Volume 3} (OpenStax), Sections 4.1 and 4.2.\par}
\end{frame}

% =============================================================
\section{Functions of several variables}

\begin{frame}{Examples}
  \begin{itemize}
    \item $f(x, y) = x^2 + y^2$: \ $\R^2 \to \R$.
    \item The temperature $T(x, y, z)$ at a point of a room: \ $\R^3 \to \R$.
    \item The loss $L(\mathbf{w})$ of a machine-learning model with weights
      $\mathbf{w} \in \R^n$, often with $n$ in the millions: \ $\R^n \to \R$.
    \item Polar coordinates $(r, \theta) \mapsto (r\cos\theta, r\sin\theta)$: \ $\R^2 \to \R^2$.
  \end{itemize}
  \pause
  We write $\vx = (x_1, \dots, x_n)$ for a point of $\R^n$, and $f\colon D \to \R^m$ with a
  \textbf{domain} $D \subseteq \R^n$. For $m = 1$, $f$ is \textbf{real-valued}; for $m > 1$,
  \textbf{vector-valued}, $f = (f_1, \dots, f_m)$.
\end{frame}

\begin{frame}{Graphs and level sets}
  \begin{columns}[c]
    \begin{column}{0.5\textwidth}
      \begin{tikzpicture}
        \begin{axis}[width=\linewidth, height=5.8cm, view={35}{25},
            axis lines=center, axis line style={wosm-muted},
            ticks=none, xlabel={$x$}, ylabel={$y$}, zlabel={$z$},
            label style={font=\scriptsize, text=wosm-muted},
            xmin=-1.6, xmax=1.6, ymin=-1.6, ymax=1.6, zmin=0, zmax=2.4]
          \addplot3[mesh, draw=wosm-accent!45, domain=-1.1:1.1, y domain=-1.1:1.1,
            samples=13, samples y=13] {x^2 + y^2};
          \addplot3[wosm-brick, very thick, domain=0:360, samples=60, samples y=0]
            ({0.5*cos(x)}, {0.5*sin(x)}, {0.25});
          \addplot3[wosm-brick, very thick, domain=0:360, samples=60, samples y=0]
            ({cos(x)}, {sin(x)}, {1});
          \addplot3[wosm-brick, very thick, domain=0:360, samples=60, samples y=0]
            ({1.4*cos(x)}, {1.4*sin(x)}, {1.96});
        \end{axis}
      \end{tikzpicture}
    \end{column}
    \begin{column}{0.48\textwidth}
      \small
      \begin{itemize}
        \item The \textbf{graph} of $f\colon \R^2 \to \R$ is the surface $z = f(x, y)$ in $\R^3$.
          Here $z = x^2 + y^2$, a paraboloid.
        \item A \textbf{level set} is $\{\vx : f(\vx) = c\}$. For $x^2 + y^2$ the level sets are
          circles of radius $\sqrt c$ ({\color{wosm-brick}drawn at height $c$}).
        \item Contour lines on a map are level sets of the height.
      \end{itemize}
      For $n \ge 3$ we cannot draw the graph; level sets and slices still help.
    \end{column}
  \end{columns}
\end{frame}

% =============================================================
\section{Norms and distance}

\begin{frame}{The Euclidean norm}
  For $\vx, \vy \in \R^n$ (Linear Algebra):
  \[ \vx \cdot \vy = \sum_{i=1}^n x_i y_i, \qquad
     \norm{\vx} = \sqrt{\vx \cdot \vx} = \sqrt{x_1^2 + \dots + x_n^2}, \qquad
     \text{distance } \norm{\vx - \vy} . \]
  \pause
  \begin{itemize}
    \item \textbf{Cauchy--Schwarz}: $|\vx \cdot \vy| \le \norm{\vx}\, \norm{\vy}$.
    \item \textbf{Triangle inequality}: $\norm{\vx + \vy} \le \norm{\vx} + \norm{\vy}$, since
      \[ \norm{\vx + \vy}^2 = \norm{\vx}^2 + 2\,\vx \cdot \vy + \norm{\vy}^2
         \le \norm{\vx}^2 + 2\norm{\vx}\norm{\vy} + \norm{\vy}^2 = (\norm{\vx} + \norm{\vy})^2 . \]
    \item Hence also $\bigl|\norm{\vx} - \norm{\vy}\bigr| \le \norm{\vx - \vy}$.
  \end{itemize}
\end{frame}

\begin{frame}{Norms}
  \begin{definition}
    A \textbf{norm} on $\R^n$ is a function $\vx \mapsto \norm{\vx} \ge 0$ with
    (1) $\norm{\vx} = 0$ only for $\vx = \mathbf{0}$, \
    (2) $\norm{\lambda \vx} = |\lambda|\, \norm{\vx}$, \
    (3) $\norm{\vx + \vy} \le \norm{\vx} + \norm{\vy}$.
  \end{definition}
  \pause
  \begin{columns}[c]
    \begin{column}{0.52\textwidth}
      \small
      \begin{itemize}
        \item {\color{wosm-accent}$\norm{\vx}_2 = \sqrt{\sum x_i^2}$}: Euclidean (straight line).
        \item {\color{wosm-brick}$\norm{\vx}_1 = \sum |x_i|$}: Manhattan (along a street grid).
        \item $\norm{\vx}_\infty = \max_i |x_i|$: maximum (the moves of a chess king).
      \end{itemize}
      The ``unit balls'' $\{\vx : \norm{\vx} \le 1\}$ in $\R^2$: a disc, a diamond, a square.
    \end{column}
    \begin{column}{0.44\textwidth}
      \centering
      \begin{tikzpicture}[scale=1.15]
        \draw[wosm-muted, ->] (-1.4,0) -- (1.5,0);
        \draw[wosm-muted, ->] (0,-1.4) -- (0,1.5);
        \draw[wosm-text, thick] (-1,-1) rectangle (1,1);
        \draw[wosm-accent, thick] (0,0) circle (1);
        \draw[wosm-brick, thick] (1,0) -- (0,1) -- (-1,0) -- (0,-1) -- cycle;
        \node[font=\scriptsize, text=wosm-muted, below right] at (1,0) {$1$};
      \end{tikzpicture}
    \end{column}
  \end{columns}
\end{frame}

\begin{frame}{All norms on $\R^n$ are equivalent}
  \begin{theorem}
    For every $\vx \in \R^n$: \quad
    $\norm{\vx}_\infty \le \norm{\vx}_2 \le \norm{\vx}_1 \le n \norm{\vx}_\infty$.
  \end{theorem}
  \pause
  {\small Proof: $\max_i x_i^2 \le \sum x_i^2 \le \bigl(\sum |x_i|\bigr)^2$, and each
  $|x_i| \le \max_j |x_j|$.\par}
  \pause
  \begin{itemize}
    \item So if $\norm{\vx_k - \vx} \to 0$ in one of these norms, it does in all of them.
      In fact any two norms on $\R^n$ satisfy such inequalities (we omit the proof).
    \item Everything in this lecture --- open sets, limits, continuity --- is the same for every
      norm. We use $\norm{\cdot} = \norm{\cdot}_2$ unless we say otherwise.
    \item Consequence: $\vx_k \to \vx$ if and only if each coordinate converges,
      $x_{k,i} \to x_i$ \ (use $\norm{\cdot}_\infty$).
  \end{itemize}
  {\small In infinite-dimensional spaces (functions, Lecture 4) norms need not be equivalent:
  uniform and pointwise convergence differ.\par}
\end{frame}

% =============================================================
\section{Open and closed sets}

\begin{frame}{Balls, interior and boundary}
  \begin{columns}[c]
    \begin{column}{0.58\textwidth}
      \begin{itemize}
        \item \textbf{Open ball}: $B(\va, r) = \{\vx : \norm{\vx - \va} < r\}$.
        \item $\va$ is an \textbf{interior point} of $S$ if some ball $B(\va, r)$, $r > 0$,
          lies inside $S$.
        \item $\va$ is a \textbf{boundary point} of $S$ if every ball $B(\va, r)$ contains points
          of $S$ and points outside $S$.
      \end{itemize}
      \begin{definition}
        $S$ is \textbf{open} if all its points are interior points. \
        $S$ is \textbf{closed} if its complement $\R^n \setminus S$ is open.
      \end{definition}
    \end{column}
    \begin{column}{0.4\textwidth}
      \centering
      \begin{tikzpicture}[scale=0.95]
        \fill[wosm-accent!14] plot[smooth cycle, tension=0.8]
          coordinates {(0,0) (2,-0.3) (3,0.8) (2.4,2.2) (0.8,2.1) (-0.3,1.2)};
        \draw[wosm-accent, thick, dashed] plot[smooth cycle, tension=0.8]
          coordinates {(0,0) (2,-0.3) (3,0.8) (2.4,2.2) (0.8,2.1) (-0.3,1.2)};
        \fill[wosm-text] (1.3,1) circle (1.2pt);
        \draw[wosm-text] (1.3,1) circle (0.45);
        \node[font=\scriptsize, text=wosm-text, below] at (1.3,0.95) {$\va$};
        \fill[wosm-brick] (3,0.8) circle (1.2pt);
        \draw[wosm-brick] (3,0.8) circle (0.35);
        \node[font=\scriptsize, text=wosm-brick, right] at (3.35,0.8) {$\mathbf{b}$};
        \node[font=\scriptsize, text=wosm-accent] at (0.6,0.4) {$S$};
      \end{tikzpicture}
      \par\small\color{wosm-muted}
      $\va$ is interior; $\mathbf{b}$ is a boundary point.
    \end{column}
  \end{columns}
\end{frame}

\begin{frame}{Examples}
  \begin{itemize}
    \item An open ball is open: if $\norm{\vx - \va} < r$, then $B(\vx, \varepsilon) \subseteq B(\va, r)$
      for $\varepsilon = r - \norm{\vx - \va}$, by the triangle inequality.
    \item The closed ball $\{\vx : \norm{\vx - \va} \le r\}$ is closed. Its boundary is the sphere
      $\norm{\vx - \va} = r$.
    \item The half-plane $\{(x, y) : y > 0\}$ is open; $\{(x, y) : y \ge 0\}$ is closed.
    \item $[0, 1) \times [0, 1)$ is neither open nor closed.
    \item $\R^n$ and $\varnothing$ are both open and closed.
  \end{itemize}
  \pause
  \begin{warning}
    ``Not open'' does not mean ``closed''. Most sets are neither.
  \end{warning}
\end{frame}

\begin{frame}{Closed sets and sequences}
  \begin{theorem}
    $S$ is closed if and only if it contains the limit of every convergent sequence of points
    of $S$.
  \end{theorem}
  {\small Proof idea: if $\vx_k \in S$, $\vx_k \to \vx$ and $\vx \notin S$, then $\vx$ is in the
  open complement, so a ball around $\vx$ misses $S$ --- but the $\vx_k$ get into every ball
  around $\vx$. The converse is similar.\par}
  \pause
  \begin{itemize}
    \item $S$ is \textbf{bounded} if it lies in some ball $B(\mathbf{0}, R)$.
    \item \textbf{Bolzano--Weierstrass in $\R^n$}: every bounded sequence has a convergent
      subsequence. (Apply the one-variable theorem of Calculus I, Unit 4, to one coordinate after
      another.)
    \item Closed and bounded sets are called \textbf{compact}. They are where continuous functions
      attain their maximum and minimum (Section 5).
  \end{itemize}
\end{frame}

% =============================================================
\section{Limits}

\begin{frame}{Limits in $\R^n$}
  \begin{definition}
    Let $f\colon D \to \R^m$ and let every ball around $\va$ contain points of $D$ other than
    $\va$. Then $\lim_{\vx \to \va} f(\vx) = L$ if for every $\varepsilon > 0$ there is a
    $\delta > 0$ with
    \[ \norm{f(\vx) - L} < \varepsilon \quad\text{for all } \vx \in D
       \text{ with } 0 < \norm{\vx - \va} < \delta . \]
  \end{definition}
  \pause
  \begin{itemize}
    \item The same words as in Calculus I, with $|\cdot|$ replaced by $\norm{\cdot}$.
      The limit laws (sums, products, quotients) hold with the same proofs.
    \item A vector-valued $f$ has a limit if and only if each component $f_i$ does.
    \item New: $\vx$ may approach $\va$ along infinitely many paths. The limit must be the same
      along all of them.
  \end{itemize}
\end{frame}

\begin{frame}{Different paths, different limits}
  \begin{example}[$f(x, y) = \dfrac{xy}{x^2 + y^2}$ at $(0, 0)$]
    \begin{itemize}
      \item Along the $x$-axis ($y = 0$): \ $f(x, 0) = 0 \to 0$.
      \item Along the diagonal ($y = x$): \ $f(x, x) = \frac{x^2}{2x^2} = \frac12 \to \frac12$.
      \item Along $y = mx$: \ $f(x, mx) = \frac{m}{1 + m^2}$: a different value for each slope.
    \end{itemize}
    So $\lim_{(x,y)\to(0,0)} f(x, y)$ does not exist.
  \end{example}
  \pause
  \begin{block}{Strategy}
    To show that a limit does \emph{not} exist, find two paths with different limits.
  \end{block}
\end{frame}

\begin{frame}{Lines are not enough}
  \begin{example}[$f(x, y) = \dfrac{x^2 y}{x^4 + y^2}$ at $(0, 0)$]
    \begin{itemize}
      \item Along $x = 0$: $f = 0$. \ Along $y = mx$:
        $f(x, mx) = \frac{m x^3}{x^4 + m^2 x^2} = \frac{m x}{x^2 + m^2} \to 0$.
        Every line gives $0$.
      \item Along the parabola $y = x^2$: \ $f(x, x^2) = \frac{x^4}{2x^4} = \frac12$.
    \end{itemize}
    The limit does not exist.
  \end{example}
  \pause
  \begin{warning}
    Checking all straight lines proves nothing. To show that a limit \emph{exists}, you need an
    estimate that holds for all $\vx$ near $\va$ at once.
  \end{warning}
\end{frame}

\begin{frame}{Proving that a limit exists}
  Bound $|f(\vx) - L|$ by something that depends only on $\norm{\vx - \va}$ and tends to $0$.
  \begin{example}[$\lim_{(x,y)\to(0,0)} \frac{x^2 y}{x^2 + y^2} = 0$]
    $x^2 \le x^2 + y^2$, so $\Bigl|\frac{x^2 y}{x^2 + y^2}\Bigr| \le |y| \le \norm{(x, y)} \to 0$.
  \end{example}
  \pause
  \begin{example}[Polar coordinates]
    With $x = r\cos\theta$, $y = r\sin\theta$: \
    $\frac{x^2 y}{x^2 + y^2} = r \cos^2\theta \sin\theta$, and $|r\cos^2\theta\sin\theta| \le r \to 0$.
    \ But $\frac{xy}{x^2 + y^2} = \cos\theta\sin\theta$ depends on $\theta$: no limit.
  \end{example}
  \pause
  {\small The polar bound must not depend on $\theta$: for $\frac{x^2y}{x^4 + y^2}$ each fixed
  $\theta$ gives $0$, yet there is no limit.\par}
\end{frame}

% =============================================================
\section{Continuity}

\begin{frame}{Continuity}
  \begin{definition}
    $f\colon D \to \R^m$ is \textbf{continuous at} $\va \in D$ if $\lim_{\vx \to \va} f(\vx) = f(\va)$
    (or $\va$ is an isolated point of $D$); it is continuous on $D$ if it is continuous at every point.
  \end{definition}
  \begin{itemize}
    \item Equivalently: $\vx_k \to \va$ implies $f(\vx_k) \to f(\va)$.
    \item The coordinate functions $\vx \mapsto x_i$ are continuous. Sums, products, quotients
      (where the denominator is not $0$) and compositions of continuous functions are continuous.
    \item So polynomials and rational functions in $x_1, \dots, x_n$ are continuous where
      defined, and so are $e^{xy}$, $\sin(x^2 + y)$, and $\vx \mapsto \norm{\vx}$
      \ (as $\bigl|\norm{\vx} - \norm{\va}\bigr| \le \norm{\vx - \va}$).
  \end{itemize}
\end{frame}

\begin{frame}{Continuous in each variable is not enough}
  \begin{example}
    $f(x, y) = \frac{xy}{x^2 + y^2}$ for $(x, y) \ne (0, 0)$, \ $f(0, 0) = 0$.
    \begin{itemize}
      \item For each fixed $y$, \ $x \mapsto f(x, y)$ is continuous; for each fixed $x$,
        $y \mapsto f(x, y)$ is continuous. (At $0$: $f(x, 0) = 0$ and $f(0, y) = 0$.)
      \item But $f$ is not continuous at $(0, 0)$: along $y = x$ it equals $\frac12$.
    \end{itemize}
  \end{example}
  \pause
  Continuity in $\R^n$ is about all ways of approaching a point, not only the coordinate
  directions. We will meet the same lesson for derivatives next lecture.
\end{frame}

\begin{frame}{Continuity and open and closed sets}
  \begin{theorem}
    If $f\colon \R^n \to \R$ is continuous and $c \in \R$, then
    $\{\vx : f(\vx) < c\}$ is open, and $\{\vx : f(\vx) \le c\}$ and $\{\vx : f(\vx) = c\}$ are closed.
  \end{theorem}
  {\small Proof: if $f(\va) < c$, continuity gives a ball around $\va$ on which $f < c$. The other
  two sets have open complements $\{f > c\}$ and $\{f < c\} \cup \{f > c\}$.\par}
  \pause
  \begin{example}
    \begin{itemize}
      \item $\{(x, y) : x^2 + y^2 < 1\}$ is open; the unit circle $x^2 + y^2 = 1$ is closed.
      \item $\{\vx : x_1 \ge 0, \dots, x_n \ge 0\}$ is closed: an intersection of closed sets.
      \item The level set $\{(x, y) : x^2 - y^2 = 1\}$, a hyperbola, is closed.
    \end{itemize}
  \end{example}
\end{frame}

\begin{frame}{The extreme value theorem}
  \begin{theorem}
    A continuous function $f\colon K \to \R$ on a nonempty closed and bounded set $K \subseteq \R^n$
    attains a maximum and a minimum on $K$.
  \end{theorem}
  {\small Proof idea, as in Calculus I: take $\vx_k \in K$ with $f(\vx_k) \to \sup_K f$; a subsequence
  converges (Bolzano--Weierstrass) to a point of $K$ ($K$ is closed), and $f$ is continuous there.\par}
  \pause
  \begin{itemize}
    \item Both conditions matter: $f(x, y) = x$ has no maximum on the open disc, nor on the
      closed but unbounded half-plane $x \ge 0$.
    \item Unit 14 finds such extrema with derivatives; Unit 16 does it on curved sets such as
      circles, with Lagrange multipliers.
  \end{itemize}
\end{frame}

% =============================================================
\section{Summary and practice}

\begin{frame}{Summary}
  \begin{itemize}
    \item Functions $f\colon D \subseteq \R^n \to \R^m$; graphs and level sets.
    \item Norms $\norm{\cdot}_2$, $\norm{\cdot}_1$, $\norm{\cdot}_\infty$; on $\R^n$ all norms are
      equivalent, so convergence means convergence of each coordinate.
    \item Open sets (every point has a ball inside), closed sets (complement open, or: contain
      their limits), compact sets (closed and bounded).
    \item Limits must agree along \emph{every} path; two paths with different limits disprove a
      limit; a bound in terms of $\norm{\vx - \va}$ proves one.
    \item Continuous functions: built from coordinates by algebra and composition;
      $\{f < c\}$ is open, $\{f \le c\}$ closed; extreme values on compact sets.
  \end{itemize}
\end{frame}

\begin{frame}{Practice}
  \small
  \begin{enumerate}
    \item Describe the level sets of $f(x, y) = x^2 - y^2$ and of $g(x, y) = xy$.
    \item Show that $\norm{\vx}_1 \le \sqrt n\, \norm{\vx}_2$ for $\vx \in \R^n$.
    \item Open, closed, both or neither?
      (a) $x^2 + y^2 \le 1$ \quad (b) $0 < x < 1$ in $\R^2$ \quad (c) $y = x^2$ \quad
      (d) $x > 0, \ y \ge 0$ \quad (e) $\mathbb{Q}^2$ \quad (f) $\varnothing$
    \item Find the limit at $(0, 0)$, or show that it does not exist:
      (a) $\frac{x^2 - y^2}{x^2 + y^2}$ \quad (b) $\frac{x y^2}{x^2 + y^2}$ \quad
      (c) $\frac{\sin(x^2 + y^2)}{x^2 + y^2}$ \quad (d) $\frac{x y^2}{x^2 + y^4}$
    \item Where is $f(x, y) = \frac{x^2 + y^2}{x - y}$ continuous?
    \item Show that $\vx \mapsto \norm{\vx}$ is continuous on $\R^n$.
    \item Give a continuous function on a bounded set in $\R^2$ that has no maximum.
  \end{enumerate}
\end{frame}

\begin{frame}{Answers}
  \small
  \begin{enumerate}
    \item Hyperbolas $x^2 - y^2 = c$ (and the lines $y = \pm x$ for $c = 0$); hyperbolas $xy = c$
      (and the two axes for $c = 0$).
    \item Cauchy--Schwarz with $(|x_1|, \dots, |x_n|)$ and $(1, \dots, 1)$.
    \item (a) closed \ (b) open \ (c) closed (a level set) \ (d) neither \
      (e) neither: no ball fits inside, and irrational points are limits of rational ones \
      (f) both.
    \item (a) does not exist: $1$ along the $x$-axis, $-1$ along the $y$-axis. \
      (b) $0$: \ $|\cdot| \le |x|$. \
      (c) $1$: \ $\frac{\sin t}{t} \to 1$ with $t = x^2 + y^2 \to 0$. \
      (d) does not exist: $0$ along lines, $\frac12$ along $x = y^2$.
    \item Wherever $x \ne y$: a quotient of continuous functions.
    \item $\bigl|\norm{\vx} - \norm{\va}\bigr| \le \norm{\vx - \va}$, so take $\delta = \varepsilon$.
    \item $f(x, y) = x$ on the open unit disc: its values approach $1$ but never reach it.
  \end{enumerate}
\end{frame}

\begin{frame}{Next time}
  \textbf{Lecture 10 --- Partial derivatives and differentiability} (Unit 10)
  \begin{itemize}
    \item Partial derivatives: differentiate in one variable, freeze the others
    \item Why partial derivatives are not enough
    \item The total derivative: the best linear approximation, and the Jacobian matrix
  \end{itemize}
  \vskip1em
  {\small
  Reading: Strang \& Herman, \emph{Calculus, Volume 3}, Sections 4.3 and 4.4.\\[0.4em]
  Questions: the module channel on Discord, or office hours,
  Monday--Friday, 10:00--17:00 (Baku time, UTC+4),
  \href{mailto:ali.khudiyev@wosm.academy}{ali.khudiyev@wosm.academy}.\par}
\end{frame}

\end{document}
